\documentclass[10pt,a4paper]{amsart}
\usepackage[margin=19mm]{geometry}
\usepackage{amsmath,amssymb,mathtools}
\usepackage[scr=rsfs,cal=euler]{mathalfa}
\usepackage{xfrac}
\usepackage[unicode,hidelinks,bookmarks=false]{hyperref}
\newcommand{\ie}{i.e.\ }
\newcommand{\cf}{cf.\ }
\newcommand{\resp}{respectively\ }
\newcommand{\Subsection}{\subsection}
\providecommand{\nobreakdash}{\nobreak\hbox{-}}
\setlength{\emergencystretch}{2em}
\multlinegap0pt
\usepackage{bm}
\usepackage{bbm}
\usepackage{dsfont}
\usepackage{xspace}
\usepackage{cancel}
\usepackage{mathtools}
\usepackage{amsthm}
\makeatletter
\@ifundefined{thm}{\newtheorem{thm}{Theorem}}{}
\@ifundefined{lem}{\newtheorem{lem}[thm]{Lemma}}{}
\makeatother
\title[Lascoux series, parking functions and noncrossing partitions]{Lascoux series, parking functions and noncrossing partitions}
\author{Alice L. L. Gao, Xin-Bei Liu, Arthur L. B. Yang,, James J. Y. Zhao}
\date{Source: arXiv:2608.15100v1 (CC BY 4.0)}
\begin{document}
\maketitle

\noindent\emph{Source and licence.} Lascoux series, parking functions and noncrossing partitions. Version arXiv:2608.15100v1 (CC BY 4.0), distributed under the Creative Commons Attribution 4.0 International licence, as recorded in the arXiv licence metadata for this exact version and shown on the abstract page https://arxiv.org/abs/2608.15100v1. Full rights notice: licenses/arxiv-2608.15100-CC-BY-4.0.txt.\par\smallskip

\noindent\textit{Editorial scope.} Counted unit: the Thm whose source label is \texttt{thm-main-inequality} in the article (source0.tex lines 645--650, source label \texttt{thm-main-inequality}), delivered together with the article's own complete original proof of it (source0.tex lines 677--773, one \texttt{proof} environment of 97 lines). No mathematics is written, repaired or re-derived: statement and proof are the article's own text line for line. The proof is detached from the statement in the source: the article states the result at lines 645--650 and prints its proof at lines 677--773, and the proof environment's own header names the statement, so the association is the article's own. Required context retained uncounted: eq:defi-O\_nk (lines 235--237); eq:defi-fnk(x) (lines 239--242); eq-g-expansion (lines 604--606); eq-lc-coeff-gnmv (lines 641--643); lem-recursion (lines 653--660); eq:rec-cki (lines 656--658). Every definition, display and label of those lines is retained unchanged. Omitted from this package and not redistributed: 1--234, 238--238, 243--603, 607--640, 644--644, 651--652, 659--676, 774--1215. Nothing inside a retained excerpt is omitted or altered: every retained body is delivered exactly as the article's lines are written, so no omission receipt of any kind accompanies this package. Citations: the retained excerpts carry no citation command at all, so no bibliography entry of the article is transcribed and no reference is redistributed. Reference adaptation: none was needed and none was made; every reference inside the delivered unit resolves inside it, so no unresolved reference or citation is produced. Printed slips kept literal: no printed word, symbol, hypothesis or number of the counted statement or of its proof was altered. Layout and class: the article's own class is \texttt{article} with packages including amsmath, amssymb, latexsym, mathtools, amsthm, tikz, tabularx, booktabs, mathrsfs, enumitem, xcolor, ytableau, graphicx, algorithm; the wrapper is the lane's pinned \texttt{amsart} editorial wrapper at 10pt on A4, which restates the theorem-like environments the retained text needs and adds the article's own macro definitions that the retained text needs (none), with extra wrapper packages (bm, bbm, dsfont, xspace, cancel, mathtools). The delivered numbering is the unit's own. Assets: no retained span contains a figure environment, a graphics-inclusion command, a PDF-page inclusion, a table or a TikZ picture; no image file is copied into this package, the package declares no asset dependency and no image file is redistributed with it. Licence: version arXiv:2608.15100v1 of this article is distributed under the Creative Commons Attribution 4.0 International licence, as recorded in the arXiv licence metadata for this exact version and shown on the abstract page https://arxiv.org/abs/2608.15100v1. A full rights notice is delivered at licenses/arxiv-2608.15100-CC-BY-4.0.txt. Characters: every retained excerpt is pure ASCII, so the delivered engine, XeLaTeX with its default Latin Modern text font, reports no missing character.
\begin{align}\label{eq:defi-O_nk}
O^{\bf{v}}_{n,m}:=(nD - v_m) \cdots(nD - v_1)(nD-v_0).
\end{align}

\begin{align}\label{eq:defi-fnk(x)}
O^{\bf{v}}_{n,m} \left( \frac{1}{1 - t} \right)=
\frac{g^{\bf{v}}_{n,m}(t)}{(1 - t)^{m + 2}},
\end{align}

\begin{align}\label{eq-g-expansion}
g^{\mathbf{v}}_{n,m}(t) = \sum_{j=0}^{m+1} c^{\mathbf{v}}_{n,m,j} t^j.
\end{align}

\begin{align}\label{eq-lc-coeff-gnmv}
(c^{\mathbf{v}}_{n,m,j})^2\geq c^{\mathbf{v}}_{n,m,j-1}c^{\mathbf{v}}_{n,m,j+1}.
\end{align}

\begin{lem}\label{lem-recursion}
For given $n,m$ and $\mathbf v=(v_0,v_1,\ldots,v_m)$, let $c^{\mathbf{v}}_{n,m,k}$ be as defined by \eqref{eq-g-expansion}. 
If $v_0=0$, then for any $0\le j \le m$ we have 
\begin{align}\label{eq:rec-cki}
 c^{\mathbf{v}}_{n,m,j}=(n(m+1-j)+v_m)c^{\overline{\mathbf{v}}}_{n,m-1,j-1}+(n(j+1)-v_m)c^{\overline{\mathbf{v}}}_{n,m-1,j},
\end{align}
where $\overline{\mathbf{v}}=(v_0,v_1,\ldots,v_{m-1})$ and $c^{(v_0)}_{n,0,0}=n$.
\end{lem}

\begin{thm}\label{thm-main-inequality} 
If $v_0=0$ and $0\le v_i<n$ for each $1\leq i\leq m$, then 
    \begin{align}\label{conj-1-inq-A}
        \frac{c^{\mathbf{v}}_{n,m,0}}{c^{\mathbf{v}}_{n,m,m}}\le\frac{c^{\mathbf{v}}_{n,m,1}}{c^{\mathbf{v}}_{n,m,m-1}}\le\cdots\le\frac{c^{\mathbf{v}}_{n,m,m}}{c^{\mathbf{v}}_{n,m,0}}.
    \end{align}
\end{thm}

\begin{proof}[Proof of Theorem \ref{thm-main-inequality}.]
Use induction on $m$. Note that the $m=0$ case is vacuous. 
From \eqref{eq:defi-O_nk} and \eqref{eq:defi-fnk(x)} it follows that 
$g^{{(v_0,v_1)}}_{n,1}(t)=n(n+v_1)t+n(n-v_1)$. Since $0\leq v_1<n$, 
the assertion is clear for $m=1$.
Now let $m\geq 2$ and assume the assertion for $m-1$.
It is sufficient to show that for each $0\leq j\leq m-1$,
\begin{align*}%\label{conj-1-inq-B}
        \frac{c^{\mathbf{v}}_{n,m,j}}{c^{\mathbf{v}}_{n,m,m-j}}\le\frac{c^{\mathbf{v}}_{n,m,j+1}}{c^{\mathbf{v}}_{n,m,m-1-j}}.
    \end{align*} 
We will use the recurrence relation \eqref{eq:rec-cki}. Since $n$ and $\mathbf{v}$ will be fixed throughout the proof, the use of $c_{m,j}$ to denote $c^{\mathbf v}_{n,m,j}$ should cause no confusion. 
By Lemma \ref{lem-recursion} and the hypothesis that $0\leq v_i<n$ for each $0\leq i\leq m$,
we see that $c_{m,j}> 0$ for any $0 \le j \le m$.
Thus it remains to prove the following equivalent inequality
    \begin{align}\label{conj-1-inq-C}
      c_{m,j+1}c_{m,m-j}-  c_{m,j}c_{m,m-1-j}\ge 0.
    \end{align}
For any $0\leq i,i' \leq m-1$, let
\begin{align*}
\mathcal{D}_{m-1,i,i'}:=c_{m-1,i}c_{m-1,m-1-i'}-c_{m-1,i'}c_{m-1,m-1-i}.
 \end{align*}
The induction hypothesis
    \begin{align*}
        \frac{c_{m-1,0}}{c_{m-1,m-1}}\le\frac{c_{m-1,1}}{c_{m-1,m-2}}\le\cdots\le\frac{c_{m-1,m-1}}{c_{m-1,0}}
    \end{align*}
implies that $\mathcal{D}_{m-1,i,i'}$ is nonnegative whenever $i\geq i'$. 
By Lemma \ref{lem-recursion}, for any $0\leq j \leq m-1$ we have 
\begin{align*}
c_{m,j}&=\big(n(m+1-j)+v_m\big)c_{m-1,j-1}+\big(n(j+1)-v_m\big)c_{m-1,j},\\
c_{m,m-j}&=\big(n(j+1)+v_m\big)c_{m-1,m-1-j}+\big(n(m+1-j)-v_m\big)c_{m-1,m-j},\\
c_{m,j+1}&=\big(n(m-j)+v_m\big)c_{m-1,j}+\big(n(j+2)-v_m\big)c_{m-1,j+1},\\
c_{m,m-1-j}&=\big(n(j+2)+v_m\big)c_{m-1,m-2-j}+\big(n(m-j)+v_m\big)c_{m-1,m-1-j}.
\end{align*}
One can verify that substituting into 
the left hand side of \eqref{conj-1-inq-C} yields
\begin{align}
       c_{m,j+1}c_{m,m-j}&- c_{m,j}c_{m,m-1-j}\notag\\
       =& (nm-nj-v_m)(nm-nj+n+v_m)\mathcal{D}_{m-1,j,j-1} \notag \\ 
       &+(nj+n-v_m)(nj+2n+v_m)\mathcal{D}_{m-1,j+1,j} \notag\\ 
      &+(nj+2n-v_m)(nm-nj+n-v_m)\mathcal{D}_{m-1,j+1,j-1} \notag \\
      &+2v_mnT_{n,m,j}, \qquad \qquad \qquad ~~~\label{conj-1-inq-D}
 \end{align}
 where 
 \begin{align*}
   T_{n,m,j}&=(c_{m-1,j} c_{m-1,m-j}+c_{m-1,j+1}c_{m-1,m-1-j}-2c_{m-1,j-1}c_{m-1,m-2-j})\\
   &+(m+1)(c_{m-1,j}c_{m-1,m-1-j} -c_{m-1,j-1}c_{m-1,m-2-j}). 
 \end{align*}
In view of the nonnegativity of $\mathcal{D}_{m-1,j,j-1}$, $\mathcal{D}_{m-1,j+1,j}$ and $\mathcal{D}_{m-1,j+1,j-1}$, 
it suffices to show that 
\begin{align}\label{eq-1-sim}
c_{m-1,j} c_{m-1,m-j}+c_{m-1,j+1}c_{m-1,m-1-j}-2c_{m-1,j-1}c_{m-1,m-2-j}\ge 0
\end{align}
and
\begin{align}\label{eq-2-sim}
c_{m-1,j}c_{m-1,m-1-j} -c_{m-1,j-1}c_{m-1,m-2-j}\ge 0
\end{align}
for any $0 \leq j \leq m - 1$. We may assume that $1 \leq j \leq m - 2$ since \eqref{eq-1-sim} and \eqref{eq-2-sim} clearly hold for $j=0$ or $j=m-1$.
Note that \eqref{eq-1-sim} is equivalent to 
\begin{align}\label{eq-3-sim}
\frac{c_{m-1,j}}{c_{m-1,j-1}} \cdot \frac{c_{m-1,m-j}}{c_{m-1,m-2-j}}+\frac{c_{m-1,j+1}}{c_{m-1,j-1}}\cdot\frac{c_{m-1,m-1-j}}{c_{m-1,m-2-j}}\ge 2.
\end{align}
From the induction hypothesis, we have
\begin{align}\label{eq-3-sim-2}
\frac{c_{m-1,j}}{c_{m-1,m-1-j}}\ge \frac{c_{m-1,j-1}}{c_{m-1,m-j}} \quad \text{and} \quad 
\frac{c_{m-1,j+1}}{c_{m-1,m-2-j}}\ge\frac{c_{m-1,j-1}}{c_{m-1,m-j}},
\end{align}
which are equivalent to
$$\frac{c_{m-1,j}}{c_{m-1,j-1}}\ge \frac{c_{m-1,m-1-j}}{c_{m-1,m-j}} \quad \text{and}\quad
\frac{c_{m-1,j+1}}{c_{m-1,j-1}}\ge \frac{c_{m-1,m-2-j}}{c_{m-1,m-j}}.$$
Substituting the above two inequalities into the left-hand side of
\eqref{eq-3-sim}, we obtain
 \begin{align*}%\label{ineq-1}
     \frac{c_{m-1,j}}{c_{m-1,j-1}}\cdot\frac{c_{m-1,m-j}}{c_{m-1,m-2-j}}+&\frac{c_{m-1,j+1}}{c_{m-1,j-1}}\cdot\frac{c_{m-1,m-1-j}}{c_{m-1,m-2-j}}
     \\[5pt]
    & 
   \ge \frac{c_{m-1,m-1-j}}{c_{m-1,m-j}}\cdot\frac{c_{m-1,m-j}}{c_{m-1,m-2-j}}+\frac{c_{m-1,m-2-j}}{c_{m-1,m-j}}\cdot\frac{c_{m-1,m-1-j}}{c_{m-1,m-2-j}} \\[5pt]
     & =\frac{c_{m-1,m-1-j}}{c_{m-1,m-2-j}}+\frac{c_{m-1,m-1-j}}{c_{m-1,m-j}}\\[5pt]
     & \geq \frac{c_{m-1,m-j}}{c_{m-1,m-1-j}}+\frac{c_{m-1,m-1-j}}{c_{m-1,m-j}}\geq 2,
 \end{align*}
 as desired, where the second inequality follows from  \eqref{eq-lc-coeff-gnmv}. 
Let us proceed to prove \eqref{eq-2-sim}, which is equivalent to
\begin{align*}%\label{eq-4-sim}
\frac{c_{m-1,j-1}} {c_{m-1,m-1-j}}
\le
\frac{c_{m-1,j}}{c_{m-1,m-2-j}}.
\end{align*}
By using \eqref{eq-lc-coeff-gnmv} and \eqref{eq-3-sim-2}, one can show that
\begin{align*}
\frac{c_{m-1,j-1}} {c_{m-1,m-1-j}}
=\frac{c_{m-1,j-1}}{c_{m-1,m-j}} \cdot \frac{c_{m-1,m-j}}{c_{m-1,m-1-j}} 
\le
\frac{c_{m-1,j}}{c_{m-1,m-1-j}} \cdot \frac{c_{m-1,m-1-j}}{c_{m-1,m-2-j}}
=\frac{c_{m-1,j}}{c_{m-1,m-2-j}},
\end{align*}
as desired.
This completes the proof.
\end{proof}

\end{document}
